\section{What the quotes determine}
\label{sec:lvinfo}

Before writing an objective function, a chapter should ask whether the
problem it is about to pose has one answer.  This one does not, in three
distinct ways --- one about resolution, one about averages, one about
whole regions of the field --- and each dictates a piece of the
calibration's design.

\paragraph{The length scale.}
Everything interesting about a smile at maturity $\tau$ happens within a
few multiples of the diffusive width $\sigma\sqrt{\tau}$ of the money.
This single quantity dictates where both grids of
\cref{sec:lvsheet,sec:lvlattice} must put their points.  \emph{Vertex
placement}: strike columns are spaced by delta --- a fixed ladder of
exercise probabilities mapped through $\PhiN^{-1}$, rather than fixed
steps of strike --- so parameter density follows where optionality
lives.  \emph{Vertex coverage}: an axis sized by
the longest expiry undersamples a short one, whose entire smile can fall
between two columns --- a two-day smile at $20\%$ volatility is about
$1.5\%$ wide, an order of magnitude narrower than a one-year smile --- so
every expiry must be guaranteed several columns \emph{within its own}
range, spread evenly across it.  \emph{Lattice steps}: the same width
bounds the strike step from above, and if the \emph{quote} spacing is
finer than the lattice step, the fitted smile oscillates at quote
frequency until the lattice out-resolves it.  A grid that does not resolve
the data's length scale fits an alias, not the data; the exact placement
rules used in this chapter are collected in \cref{app:lvnumerics}.

\paragraph{An averaging fact, and the floor.}
At the money, implied and local variance are related as average to
integrand.  If the field near the money depends only on time,
$v(\tau,y)\approx v(\tau,1)$, the model is Black--Scholes with a
time-dependent volatility and \emph{exactly}
$w(0,\tau)=\int_0^{\tau}v(s,1)\,\dd s$; in general this holds to leading
order near the money.  An average can only be matched by an integrand that
crosses it: whenever the ATM term structure is \emph{increasing}, early
local variance must sit strictly \emph{below} the shortest implied
variance.  A parameter floor set carelessly --- say at a round $5\%$
volatility --- therefore makes some low-volatility short-dated term
structures literally unfittable, with the optimizer pinned to the box and
the misfit looking like a mystery.  The box of \cref{prop:lvnodal} must be
generous and \emph{adaptive}: a cap a fixed multiple of the largest
implied volatility in view, a floor a fraction of the \emph{smallest ATM}
implied --- keyed to ATM quotes deliberately, so that one noisy deep-wing
print cannot lower the floor where it does real stabilizing work.

\paragraph{Rows the quotes cannot see.}
The quotes at the first expiry $\tau_1$ pin --- through the averaging fact
--- only the \emph{integral} of local variance over $[0,\tau_1]$.  Any
vertex row strictly inside that interval is individually unidentified: the
optimizer can ring one row up and the next down at no cost in fit, and
under noise it will (this is precisely how the rippled fit of
\cref{fig:lvwrongway} spent its freedom at the shortest expiry, wandering
\MacLvWrongEdgePts\ volatility points at a span edge held by a single
quote).  The cure is not more data but a declared prior: tie each sub-front
row to the first identified row --- the \emph{front tie} of
\cref{sec:lvcalib} --- turning an unidentified subspace into a stated
constant continuation.

\paragraph{Two sheets, one quote set.}
The three paragraphs above are instances of one geometry, and it is better
seen than told.  A fit like this chapter's carries a couple of hundred
nodal parameters, and even when the quotes outnumber them several-fold
--- nearly five to one on the frozen SPY chain of \cref{sec:lvexamples},
between two and three on NVDA --- the quotes cluster near the money of a
handful of expiries, and a vanilla constrains local variance only through
a smoothing integral along diffusion paths.  Whole regions of the sheet
(deep wings, sub-front rows, times beyond the last expiry) are therefore
weakly determined no matter how the counting goes.  \Cref{fig:lvidentify} makes it concrete on the synthetic
running example: pushing the sheet's unquoted deep-put column (at
$y=\MacLvIdentColY$, below every quoted strike) down by
\MacLvIdentMovePts\ full volatility points changes the reprice of the
entire quote set by at most \MacLvIdentMaxDiffBp\ bp ---
\MacLvIdentRmsDiffBp\ bp rms --- concentrated in the longest expiry's
lowest strikes, the only quotes whose diffusion reaches the moved column
at all.  A thousand basis points of surface, twenty of quotes.

\begin{figure}[htbp]
\centering
\paperfig{\textwidth}{fig_lv_identify.pdf}
\caption{What sparse quotes do not see.  (a)~The calibrated synthetic
sheet at $\tau=0.5$ (solid) and the same sheet with its unquoted deep-put
column pushed down by \MacLvIdentMovePts\ volatility points (dashed); the
shaded band is the quoted region, where the two are indistinguishable.
(b)~The absolute reprice change at every quote of every expiry: at most
\MacLvIdentMaxDiffBp\ bp, and only where the longest expiry's lowest
strikes feel the moved column.  Low quote error is \emph{not} surface
recovery.}
\label{fig:lvidentify}
\end{figure}

The consequence is philosophical and operational at once.  What a
calibration returns is one \emph{representative} of an equivalence class
of sheets the data cannot distinguish, and every device in the next
section's objective is a declared choice of representative: the roughness
term selects the smoothest member near a reference, the box clips the
unidentified extremes, the front tie pins the rows below the first expiry,
and the seed decides which basin a non-convex search lands in.  None of
this is hidden tuning; it is the ``convention'' column of the book's
thesis, operating inside a calibration.  The worked example of
\cref{sec:lvexamples} accordingly reports quote error and surface error as
two separate numbers.
